\subsection{多重模}

设 A、B 为两个环，并在集合 E 上考虑具有相同加法法则且算子环分别为 A 和 B 的两个左模结构；设 \(\mathcal{E}\) 为加法群 E 的自同态环，且对所有 \(\alpha \in A\) （相应地 \(\beta \in B\) ) 令 \(h_{\alpha}\) （相应地 \(h_{\beta}^{\prime}\) ) 表示元素 \(x \mapsto \alpha x\) （相应地 \(x \mapsto \beta x\) ），它们属于 \(\mathcal{E}\) 。显然，以下三个性质是等价的：(a) \(h_{\alpha} \circ h_{\beta}^{\prime} = h_{\beta}^{\prime} \circ h_{\alpha}\) 对所有 \(\alpha\) 和 \(\beta\) 成立；(b) A 在同态 \(a \mapsto h_{\alpha}\) 下的像包含于 \(\operatorname{Hom}_{\mathbf{B}}(\mathbf{E}, \mathbf{E})\) ；(c) B 在同态 \(\beta \mapsto h_{\beta}^{\prime}\) 下的像包含于 \(\operatorname{Hom}_{\mathbf{A}}(\mathbf{E}, \mathbf{E})\) 。当所涉及的 A-模（相应地 B-模）结构是右模结构时，在 (b)（相应地 (c)）中必须将环 A（相应地 B）替换为 \(\mathbf{A}^{0}\) （相应地 \(\mathbf{B}^{0}\) ）。上述性质可以表述为：定义在 E 上的两个（左或右）模结构是相容的。

定义13. 设 \((\mathbf{A}_{\lambda})_{\lambda \in \mathbf{L}}, (\mathbf{B}_{\mu})_{\mu \in \mathbf{M}}\) 为两个环族；一个 \(((\mathbf{A}_{\lambda}), (\mathbf{B}_{\mu}))\) -多重模（或关于环族 \((\mathbf{A}_{\lambda})_{\lambda \in \mathbf{L}}, (\mathbf{B}_{\mu})_{\mu \in \mathbf{M}})\) 的多重模）是一个集合 E，对于每个 \(\lambda \in \mathbf{L}\) ，它具有一个左 \(\mathbf{A}_{\lambda}\) -模结构，并且对于每个 \(\mu \in \mathbf{M}\) ，具有一个右 \(\mathbf{B}_{\mu}\) -模结构，所有这些模结构彼此相容。

当族 \((\mathbf{B}_{\mu})\) （相应地 \((\mathbf{A}_{\lambda})\) ）为空时，E 称为左（相应地，右）多重模。当 \(\operatorname{Card}(\mathbf{L}) + \operatorname{Card}(\mathbf{M}) = 2\) 时，我们称“双模”而非“多重模”；此时通常方便将双模视为（这总可以通过将算子环替换为其反环来实现，参见第1号）关于环 A 具有左模结构、关于环 B 具有右模结构，而运算律的可交换性则由关系式

\begin{enumerate}
\def\labelenumi{(\arabic{enumi})}
\setcounter{enumi}{48}
\tightlist
\item
\end{enumerate}

\[
\alpha (x \beta) = (\alpha x) \beta \quad \text { for } x \in \mathbf {E}, \alpha \in \mathbf {A}, \beta \in \mathbf {B}.
\]

然后也说 \(\mathbf{E}\) 是一个(A, B)-双模。

集合E上的两个多重模结构称为相容的，如果定义这一个或那一个多重模结构的所有E上的模结构彼此相容。

例。(1) 在环A上， \(A_{s}\) 和 \(A_{d}\) 的模结构是相容的，因此A可以被典范地视为一个(A, A)-双模。

\begin{enumerate}
\def\labelenumi{(\arabic{enumi})}
\setcounter{enumi}{1}
\tightlist
\item
  一个左A-模E在环 \(\operatorname{End}_{\mathbf{A}}(\mathbf{E})\) 上具有典范的左模结构，并且A-模与 \(\operatorname{End}_{\mathbf{A}}(\mathbf{E})\) -模结构在E上是相容的。
\end{enumerate}

显然，当E是环的两族 \((\mathbf{A}_{\lambda})_{\lambda\in L},(\mathbf{B}_{\mu})_{\mu\in M}\) 上的多重模时，E也是任意两个子族 \((\mathbf{A}_{\lambda})_{\lambda\in L},(\mathbf{B}_{\mu})_{\mu\in M'}\) 上的多重模，其中 \(A_{\lambda}\) -模和 \(B_{\mu}\) -模结构（ \(\lambda\in L'\) 和 \(\mu\in M'\) ）正是最初给定的那些。

由于多重模是带算子的交换群的特别例子，第2至10号的结果（参见第10号注记）可以应用于它们；特别地，如果E, F是两个 \(((A_{\lambda}), (B_{\mu}))\) -多重模，一个同态 \(u: E \to F\) 是一个映射，它作为 \(A_{\lambda}\) -同态对所有 \(\lambda \in L\) 成立，并作为 \(B_{\mu}\) -同态对所有 \(\mu \in M\) 成立。一个 \(((A_{\lambda}), (B_{\mu}))\) -多重模的稳定子群是 \(((A_{\lambda}), (B_{\mu}))\) -多重模（称为子多重模），由这样的子群得到的商群也是如此（称为商多重模）；对于乘积和直和也有类似结论。

设 E 为一个 \(((A_{\lambda}), (B_{\mu}))\) -多重模，且对每个 \(\lambda \in L\) （相应地 \(\mu \in M\) ) 令 \(\phi_{\lambda}: A_{\lambda}' \to A_{\lambda}\) （相应地 \(\psi_{\mu}: B_{\mu}' \to B_{\mu}\) ) 为环同态；显然，诸 \(A_{\lambda}'\) -模结构（与诸 \(\phi_{\lambda}\) 以及 E 上给定的 \(A_{\lambda}\) -模结构相关联者）与诸 \(B_{\mu}'\) -模结构（与诸 \(\psi_{\mu}\) 以及 E 上给定的 \(B_{\lambda}\) -模结构相关联者）彼此相容，从而在 E 上定义一个 \(((A_{\lambda}')\) , \((B_{\mu}')\) )-多重模结构，称其为与给定的 \(((A_{\lambda})\) , \((B_{\mu})\) )-多重模结构以及诸 \(\phi_{\lambda}\) 和 \(\psi_{\mu}\) 相关联.

若 E、F 是两个 \(((A_{\lambda}), (B_{\mu}))\) -多重模，则 E 到 F 的同态加法群记为 \(\operatorname{Hom}_{(A_{\lambda}), (B_{\mu})}(E, F)\) （或简记为 \(\operatorname{Hom}(E, F)\) ）。第2号的公式 (6) 至 (8) 显然对 \(((A_{\lambda}), (B_{\mu}))\) -多重模的同态成立，并且特别地， \(\operatorname{Hom}(E, E) = \operatorname{End}(E)\) 具有环结构；

此外， \(\operatorname{Hom}(\mathbf{E},\mathbf{F})\) 具有典范的左 \(\operatorname{End}(\mathbf{F})\) -模结构和右 \(\operatorname{End}(\mathbf{E})\) -模结构，这两种结构是相容的；换言之， \(\operatorname{Hom}(\mathbf{E},\mathbf{F})\) 具有典范的(End(F), End(E))-双模结构。

假设现在E具有一个多重模结构，其左（相应地，右）算子环一方面为诸 \(A_{\lambda}\) （ \(\lambda \in L\) ）（相应地，诸 \(B_{\mu}\) ， \(\mu \in M\) ），另一方面为另一族的环 \((\mathbf{A}_{\lambda'}')_{\lambda' \in L'}\) （相应地 \((\mathbf{B}_{\mu'}')_{\mu' \in M'}\) ）。类似地假设F具有一个多重模结构，其左（相应地，右）算子环一方面为诸 \(A_{\lambda}\) （ \(\lambda \in L\) ）（相应地，诸 \(B_{\mu}\) ， \(\mu \in M\) ），另一方面为另一族的环 \((\mathbf{A}_{\lambda''})_{\lambda'' \in L''}\) （相应地 \((\mathbf{B}_{\mu''})_{\mu'' \in M''}\) ）；为简略起见，我们将称E为一个 \(((\mathbf{A}_{\lambda}), (\mathbf{A}_{\lambda'}'); (\mathbf{B}_{\mu}), (\mathbf{B}_{\mu'})\) -多重模，且F为 \(((\mathbf{A}_{\lambda}), (\mathbf{A}_{\lambda''}); (\mathbf{B}_{\mu}), (\mathbf{B}_{\mu''}))\) -多重模。将E和F视为 \(((\mathbf{A}_{\lambda}), (\mathbf{B}_{\mu}))\) -多重模，也就是把算子限制到子族 \((\mathbf{A}_{\lambda})\) 和 \((\mathbf{B}_{\mu})\) 上。根据本号开头所述，E和F上给定的多重模结构典范地定义了环同态

\[
\mathrm{A} _ {\lambda^ {\prime}} ^ {\prime} \rightarrow \operatorname{End} _ {(\mathrm{A} _ {\lambda}), (\mathrm{B} _ {\mu})} (\mathrm{E}), \quad \mathrm{B} _ {\mu^ {\prime}} ^ {\prime 0} \rightarrow \operatorname{End} _ {(\mathrm{A} _ {\lambda}), (\mathrm{B} _ {\mu})} (\mathrm{E}),
\]

\[
\mathrm{A} _ {\lambda^ {\prime \prime}} ^ {\prime \prime} \rightarrow \operatorname{End} _ {(\mathrm{A} _ {\lambda}), (\mathrm{B} _ {\mu})} (\mathrm{F}), \quad \mathrm{B} _ {\mu^ {\prime \prime}} ^ {\prime \prime} \rightarrow \operatorname{End} _ {(\mathrm{A} _ {\lambda}), (\mathrm{B} _ {\mu})} (\mathrm{F});
\]

此外， \(\operatorname{End}_{(\mathbf{A}_{\lambda}), (\mathbf{B}_{\mu})}(\mathbf{E})\) （相应地 \(\operatorname{End}_{(\mathbf{A}_{\lambda}), (\mathbf{B}_{\mu})}(\mathbf{F})\) ) 的两个元素，若分别是取自 \(\mathbf{A}_{\lambda'}'\) 或 \(\mathbf{B}_{\mu'}^{0}\) （相应地， \(\mathbf{A}_{\lambda''}\) 或 \(\mathbf{B}_{\mu''}^{0}\) ) 的两个不同环中元素的像，则是可交换的；由此可知，上述同态在 \(\operatorname{Hom}_{(\mathbf{A}_{\lambda}), (\mathbf{B}_{\mu})}(\mathbf{E}, \mathbf{F})\) 上定义了一个多重模结构，其左算子环为诸 \(\mathbf{A}_{\lambda''}\) ( \(\lambda'' \in \mathbf{L}''\) ) 和诸 \(\mathbf{B}_{\mu'}'\) ( \(\mu' \in \mathbf{M}'\) )，其右算子环为诸 \(\mathbf{A}_{\lambda'}'\) ( \(\lambda' \in \mathbf{L}'\) ) 和诸 \(\mathbf{B}_{\mu''}^{0}(\mu'' \in \mathbf{M}'')\) .

如果现在 \(\mathbf{E}'\) 是一个 \(((\mathbf{A}_{\lambda}), (\mathbf{A}_{\lambda'}^{\prime}) ; (\mathbf{B}_{\mu}), (\mathbf{B}_{\mu'}^{\prime}))\) -多重模，且 \(\mathbf{F}'\) 是一个 \(((\mathbf{A}_{\lambda}), (\mathbf{A}_{\lambda''}^{\prime}) ; (\mathbf{B}_{\mu}), (\mathbf{B}_{\mu''}^{\prime}))\) -多重模，则 \(\mathrm{Hom}_{(\mathbf{A}_{\lambda}), (\mathbf{B}_{\mu})}(\mathbf{E}', \mathbf{F}')\) 是一个

\[
\left(\left(\mathrm{A} _ {\lambda^ {\prime \prime}} ^ {\prime \prime}\right), \left(\mathrm{B} _ {\mu^ {\prime}} ^ {\prime}\right): \left(\mathrm{A} _ {\lambda^ {\prime}} ^ {\prime}\right), \left(\mathrm{B} _ {\mu^ {\prime \prime}} ^ {\prime \prime}\right)\right) - \text { multimodule };
\]

中的多重模，且若 \(u: \mathbf{E}' \to \mathbf{E}\) , \(v: \mathbf{F} \to \mathbf{F}'\) 是多重模同态，

\[
\operatorname{Hom} (u, v): \operatorname{Hom} _ {\left(\mathrm{A} _ {\lambda}\right), \left(\mathrm{B} _ {\mu}\right)} (\mathrm{E}, \mathrm{F}) \rightarrow \operatorname{Hom} _ {\left(\mathrm{A} _ {\lambda}\right), \left(\mathrm{B} _ {\mu}\right)} \left(\mathrm{E} ^ {\prime}, \mathrm{F} ^ {\prime}\right)
\]

定义如第2号所述，且是一个多重模同态。

注记. (1) 设 F 是一个 A-模，C 是环 A 的中心；由于中心位似与所有位似可交换，F 具有一个双模结构，其左算子环为 A 和 C。若 E 是另一个 A-模，则 \(\operatorname{Hom}_{\mathbf{A}}(\mathbf{E}, \mathbf{F})\) 因此具有典范的C-模结构（其中，对于 \(f \in \operatorname{Hom}_{\mathbf{A}}(\mathbf{E}, \mathbf{F})\) 和 \(\gamma \in C\) , \(\gamma f\) 是同态 \(x \mapsto \gamma f(x)\) ）；如果 \(E'\) , \(F'\) 是两个A-模， \(u: E' \to E\) , \(v: F \to F'\) 是两个A-同态，则映射 \(\operatorname{Hom}(u, v)\) 是C-线性的。

\begin{enumerate}
\def\labelenumi{(\arabic{enumi})}
\setcounter{enumi}{1}
\tightlist
\item
  设 E 为左 A-模；由于 A 具有典范的 (A, A)-双模结构，对任意指标集 T，直和 \(\mathbf{A}^{(\mathrm{T})}\) 亦然；根据上述内容， \(\mathrm{Hom}_{\mathbf{A}}(\mathbf{A}_{s}^{\mathrm{(T)}},\mathbf{E})\) 具有由 \(\mathbf{A}_{s}^{\mathrm{(T)}}\) 上的右A-模结构导出的典范左A-模结构：对于 \(f\in\mathrm{Hom}_{\mathbf{A}}(\mathbf{A}_{s}^{\mathrm{(T)}},\mathbf{E})\) 和 \(\alpha\in\mathbf{A},\alpha f\) 是线性映射 \(x\mapsto f(x\alpha)\) 。第11号命题17的推论2定义了一个典范映射 \(j_{\mathbf{E},\mathbf{T}}\) ，它把积模 \(\mathbf{E}^{\mathrm{T}}\) 映到 \(\mathrm{Hom}_{\mathbf{A}}(\mathbf{A}_{s}^{\mathrm{(T)}},\mathbf{E})\) ，其中 \(j_{\mathbf{E},\mathbf{T}}\) 把族 \((x_{t})_{t\in\mathbf{T}}\) 映为线性映射 \(f:\mathbf{A}_{s}^{\mathrm{(T)}}\to\mathbf{E}\) ，使得 \(f(e_{t})=x_{t}\) 对所有 \(t\in\mathbf{T}\) 成立（其中 \((e_{t})\) 是 \(\mathbf{A}_{s}^{\mathrm{(T)}}\) 的典范基）；已知（同上引文） \(j_{\mathbf{E},\mathbf{T}}\) 是双射的，并且由上文给出的 \(\mathrm{Hom}_{\mathbf{A}}(\mathbf{A}_{s}^{\mathrm{(T)}},\mathbf{E})\) 上A-模结构的定义可知， \(j_{\mathbf{E},\mathbf{T}}\) 是A线性的。最后，如果 \(u:\mathbf{E}\to\mathbf{F}\) 是一个A-模同态，则图表
\end{enumerate}

\[
\begin{array}{c} \mathrm {E^ {T}} \xrightarrow {j _ {\mathrm{E,T}}} \operatorname{Hom} _ {\mathrm{A}} (\mathrm{A} _ {s} ^ {(\mathrm{T})}, \mathrm{E}) \\ u ^ {\mathrm{T}} \Biggl \downarrow \qquad \qquad \qquad \qquad \qquad \qquad \qquad \qquad \qquad \qquad \qquad \qquad \qquad \qquad \qquad \qquad \qquad \qquad \qquad \qquad \qquad \qquad \qquad \qquad \qquad \qquad \qquad \text {Hom(1,u)} \\ \mathrm {F^ {T}} \xrightarrow {j _ {\mathrm{F,T}}} \operatorname{Hom} _ {\mathrm{A}} (\mathrm{A} _ {s} ^ {(\mathrm{T})}, \mathrm{F}) \end{array}\tag{50}
\]

是交换的。

注意，当T仅包含一个元素时， \(j_{E}:E \to \operatorname{Hom}_{A}(A_{s}, E)\) 正是第2号例1中所定义的映射 \(x \mapsto \theta_{x}\) 。

